% Agent 17: Pages 404--407 (PDF pages 32--33)
% Sections: 4.8 Accretion onto a Moving Hole,
%           4.9 Optical Appearance of Hole: Summary

%% ============================================================
%% Section 4.8: Accretion onto a Moving Hole
%% ============================================================

\subsection{Accretion onto a Moving Hole}
\label{sec:4.8}

Turn attention now from a hole at rest in the interstellar medium, to a hole that
moves with a speed much larger than the sound speed, $u_\infty \gg a_\infty$. Examine the
resulting accretion in the rest frame of the hole. At radii $r \gg 2GM/u_\infty^2$, the kinetic
energy per unit mass, $\frac{1}{2}u_\infty^2$, is far larger than the potential energy, $GM/r$; so the
pull of the hole has no effect on the gas flow. At $r \sim 2GM/u_\infty^2$, the pull of the
hole becomes significant. Because the flow is supersonic the gas will respond to
that pull in the same manner as would noninteracting particles; its pressure
cannot have an influence. (One can see this, for example, in the Euler equation

\begin{equation}
\text{const.} = \tfrac{1}{2}v^2 + \frac{a^2}{\Gamma - 1} - \frac{GM}{r}\,,
\tag{4.8.1}
\end{equation}

where the speed-of-sound (enthalpy) term---which accounts for pressure effects---is
negligible compared to the kinetic-energy term.) Thus, the flow lines will be
identical to the trajectories of test particles: they will be hyperbolae
(Fig.~4.8.1a).

After passing around the hole, the trajectories of test particles intersect
(particle collisions; dotted line of Figure~4.8.1a). In the gas-dynamic case such
intersection of flow lines is prevented by rapidly mounting pressure. Consequently,
a shock front must develop around the black hole (Fig.~4.8.1b).
Outside the shock front the flow lines will be hyperbolic test-particle trajectories.
Behind the shock front, they will not. The shock will be located roughly where
potential energy equals kinetic energy, so its characteristic size $l_s$ as shown in
Figure~4.8.1b will be
\begin{equation}
l_s \sim GM/u_\infty^2\,.
\tag{4.8.2}
\end{equation}

In the shock, the gas will lose most of its velocity perpendicular to the shock
front [``strong shock'' in terminology of \S\,2.7; ``strong'' because $u_\infty/a_\infty ={}$
(speed of gas)/(speed of sound) $\gg 1$ on front side]; but it will retain all of its
velocity parallel to the front (denote this by $u_\|$). For a given gas element, if the
remaining kinetic energy behind the front greatly exceeds the potential energy,
$\frac{1}{2}u_\|^2 \gg GM/r$, then the gas element will escape the pull of the hole. If
$\frac{1}{2}u_\|^2 \lesssim GM/r$, then it cannot escape. The dividing line between escape and
capture is a dotted line in Figure~4.8.1b; and the corresponding impact parameter
is labelled $b_\text{capture}$.

We can calculate $b_\text{capture}$ in order of magnitude by idealizing the shock as
confined to the thin dotted region of Figure~4.8.1a (Hoyle and Lyttleton~1939)\cite{HoyleLyttleton1939}.
The value of $b_\text{capture}$ will correspond to an orbit for which
\begin{equation}
\tfrac{1}{2}v_\perp^2 = GM/x \quad \text{at}\quad y = 0.
\end{equation}

Straightforward examination of Kepler orbits reveals that
\begin{equation}
b_\text{capture} = 2GM/u_\infty^2\,.
\tag{4.8.3}
\end{equation}

Since the mass flux in the gas at large radii is $\rho_\infty u_\infty$, the accretion rate is
\begin{equation}
\dot{M}_0 = (\pi b_\text{capture}^2)\,\rho_\infty\, u_\infty
= 4\pi G^2 M^2 \rho_\infty / u_\infty^3\,.
\tag{4.8.4}
\end{equation}

\begin{figure}[htbp]
\centering
% \includegraphics[width=\columnwidth]{fig_4_8_1}
\fbox{\parbox{0.9\columnwidth}{\centering [Figure 4.8.1 placeholder]\\[6pt]
(a)~Trajectories of test particles (and of supersonic gas)
relative to a black hole, as viewed in the rest frame of the hole.\\[4pt]
(b)~The flow lines of supersonic gas
relative to a black hole, as viewed in the rest frame of the hole.
The trajectories and flow lines are virtually the same, until
the gas encounters the shock front.}}
\caption{\textit{Figure~4.8.1.} The trajectories of test particles~(a), and the flow lines
of supersonic gas~(b) relative to a black hole, as viewed in the rest frame of the hole.
The trajectories and flow lines are virtually the same, until the gas encounters the
shock front.}
\label{fig:4.8.1}
\end{figure}

More careful calculations give accretion rates which differ from this by
multiplicative factors of ${\sim}\,0.5$ to $1.0$ (Bondi and Hoyle~1944)\cite{BondiHoyle1944}.

Notice that the accretion rate~(4.8.4) for the supersonic case, and the rate
(4.2.10) for a hole at rest are identical except for the replacement of sound
speed $a_\infty$ by speed of flow $u_\infty$, and except for a numerical coefficient of order
unity. In the intermediate case of $u_\infty \sim a_\infty$, one can use the hybrid formula
\begin{equation}
\dot{M}_0 \simeq \frac{4\pi G^2 M^2 \rho_\infty}{(a_\infty^2 + u_\infty^2)^{3/2}}
\tag{4.8.5}
\end{equation}

(Bondi~1952, Hunt~1971)\cite{Bondi1952,Hunt1971}.
\begin{equation}
\dot{M}_0 \simeq \left(1 \times 10^{11}~\frac{\text{g}}{\text{sec}}\right)
\frac{(M/M_\odot)^2\;(\rho_\infty/10^{-24}~\text{g~cm}^{-3})}
{\bigl[(u_\infty/10~\text{km~sec}^{-1})^2 + (a_\infty/10~\text{km~sec}^{-1})^2\bigr]^{3/2}}\,.
\notag
\end{equation}

[Salpeter~(1964)\cite{Salpeter1964} gave the first application of these formulas to black holes; all
previous work dealt with normal stars.]

It is worth noting that a significant fraction of the kinetic energy released in
the shock
\begin{equation}
\frac{dE}{dt} \simeq \tfrac{1}{2}\dot{M}_0\, u_\infty^2
\simeq \left(10^{24}~\text{ergs~sec}^{-1}\right)
\left(\frac{M}{M_\odot}\right)^{\!2}
\left(\frac{\rho_\infty}{10^{-24}~\text{g~cm}^{-3}}\right)
\left(\frac{u_\infty}{10~\text{km~sec}^{-1}}\right)^{-1/2}
\tag{4.8.6}
\end{equation}
may go into excitation of unionized atoms and thence into light as the atoms
decay, thus producing a shock front that glows (albeit weakly). See, e.g.,
Pikel'ner~(1961), and \S\S\,9 and~10 of Kaplan~(1966)\cite{Pikelner1961,Kaplan1966}.

The motion of the gas behind the shock is practically unknown, particularly
when one takes account of magnetic fields. Many different types of instabilities
may occur, both here and in the case of a stationary black hole. For example at
radii $r < l_s$, where the inflow has become supersonic again, turbulence may
create shock waves which convert kinetic energy into heat (Bisnovatyi-Kogan and
Sunyaev~1971)\cite{BisnovatyiKoganSunyaev1971}. A variety of plasma instabilities may also arise. And the
reconnection of magnetic field lines will surely be important. One can only guess that
the overall, time-averaged equipartition model of
Shvartsman~(\S\,4.4)\cite{Shvartsman1971} will resemble the equipartition model of
Shvartsman~(\S\,4.4).

One of the factors that may be important near the gravitational radius, where
most of the luminosity is produced, is angular momentum in the direction of motion
of the hole. Suppose that the accreting gas has nonzero angular momentum per
unit mass. If the angular momentum per unit mass, $\mathscr{L}$, exceeds
$r_g c$, then centrifugal forces will become important before the infalling gas reaches
the horizon; the gas will be thrown into circulating orbits; and only after viscous
stresses have transported away the excess angular momentum will the gas fall
down the hole. (See \S\,5 for a situation similar to this.) The resulting viscous
heating and lengthened time spent outside the hole, as well as the changed flow
pattern, may affect the outpouring significantly. Moreover, when the
hole moves from a region with angular momentum of gas in one direction to a
region with angular momentum in another, the gas flow may become particularly
violent for a while near the horizon; and a strong flare of luminosity may be
produced (Salpeter~1964; Shvartsman~1971)\cite{Salpeter1964,Shvartsman1971}. Almost nothing is known today
about these issues.

Is the specific angular momentum of the accreting gas likely to exceed $r_g c$?
Interstellar gas is accreted from a region of diameter
\begin{equation}
2b_\text{capture} = \frac{2r_g}{u_\infty^2/c^2}
= (3 \times 10^{14}~\text{cm})
\left(\frac{M}{M_\odot}\right)
\left(\frac{u_\infty}{10~\text{km~sec}^{-1}}\right)^{-2}\,.
\tag{4.8.7}
\end{equation}

(For $u_\infty < a_\infty$, we must replace $u_\infty$ by $a_\infty$.) The interstellar gas is turbulent with
the turbulent velocity $v_\text{turb}$ on scales $l$ given roughly by
\begin{equation}
v_\text{turb} = \left(10^5~\frac{\text{cm}}{\text{sec}}\right)
\left(\frac{l}{3 \times 10^{20}~\text{cm}}\right)^{q}\,,
\tag{4.8.8}
\end{equation}
where $q$ is a coefficient believed to lie between $1/2$ and~$1$ (Kaplan and Pikel'ner
1970)\cite{KaplanPikelner1970}, and the coefficient $10^6$~cm/sec is taken from astronomical observations.
Let us take $q = 0.75$ as a best guess. Then the specific angular momentum
associated with this turbulence, for a gas element of size $l$, is
\begin{equation}
\mathscr{L}_\text{turb} \simeq v_\text{turb}\, l
\simeq \left(3 \times 10^{26}~\frac{\text{cm}^2}{\text{sec}}\right)
\left(\frac{l}{3 \times 10^{20}~\text{cm}}\right)^{1.75}\,.
\tag{4.8.9}
\end{equation}

Consequently, the specific angular momentum of the gas that accretes from a
region of size $l = 2b_\text{capture}$ should be
\begin{equation}
\frac{\mathscr{L}_\text{turb}}{r_g c}
\simeq 1 \times
\left(\frac{M}{M_\odot}\right)^{3/4}
\left(\frac{u_\infty}{10~\text{km~sec}^{-1}}\right)^{-7/2}\,.
\tag{4.8.10}
\end{equation}

Evidently, the angular momentum may be important in some cases and may not
be in others.

%% ============================================================
%% Section 4.9: Optical Appearance of Hole: Summary
%% ============================================================

\subsection{Optical Appearance of Hole: Summary}
\label{sec:4.9}

In summary, a black hole alone in the interstellar medium may be expected to
emit ${\sim}\,10^{29}$ to $10^{35}$~ergs/sec of synchrotron radiation, concentrated in the
optical part of the spectrum (${\sim}\,10^{14}$~Hz to $10^{15}$~Hz). The output may
fluctuate on timescales of ${\sim}\,10^{-2}$ to $10^{-4}$ seconds due to instabilities, and to
orbital motion of hot spots about the hole (\S\S\,4.4, 4.7, 4.8). There may also
be strong flares when the hole moves out of one turbulent cell of interstellar
gas into another. The time interval between flares would be about
\begin{equation}
\Delta t_\text{flares}
\simeq 2b_\text{capture}/u_\infty
\simeq (10~\text{years})
\left(\frac{M}{M_\odot}\right)
\left(\frac{u_\infty}{10~\text{km~sec}^{-1}}\right)^{-3}\,.
\tag{4.9.1}
\end{equation}

It will probably be difficult observationally to distinguish accreting, isolated
black holes from accreting, old neutron stars without magnetic fields. The only
distinguishing feature will be the emission from the surface of the neutron star.
